\documentclass[11pt,a4paper]{ltjsarticle}
\usepackage[margin=25mm]{geometry}
\usepackage{amsmath,amssymb,amsthm,mathtools,booktabs,array}
\usepackage[haranoaji]{luatexja-preset}
\usepackage[unicode,hidelinks]{hyperref}
\usepackage{xurl}
\newtheorem{theorem}{定理}[section]
\newtheorem{proposition}[theorem]{命題}
\newtheorem{lemma}[theorem]{補題}
\newtheorem{corollary}[theorem]{系}
\theoremstyle{definition}
\newtheorem{definition}[theorem]{定義}
\renewcommand{\proofname}{証明}
\newcommand{\Aut}{\operatorname{Aut}}
\newcommand{\Ext}{\operatorname{Ext}}
\newcommand{\Tor}{\operatorname{Tor}}
\newcommand{\Hom}{\operatorname{Hom}}
\newcommand{\GL}{\operatorname{GL}}
\newcommand{\gr}{\operatorname{gr}}
\newcommand{\rk}{\operatorname{rank}}
\newcommand{\GKdim}{\operatorname{GKdim}}
\newcommand{\ZZ}{\mathbb Z}
\newcommand{\Gclass}{\mathcal G}
\title{最小生成元次数$(1,1,2,3)$の四次元AS正則代数\\
\large noetherian・所定Hilbert級数の下でのOre表示と次数付き同型}
\author{}
\date{2026年11月18日\quad 第1版}
\begin{document}
\hypersetup{pageanchor=false}
\begin{titlepage}
\maketitle
\thispagestyle{empty}
\begin{abstract}
複素数体上の連結局所有限な次数付き代数で、左右noetherian、
AS正則次元4、最小生成元次数$(1,1,2,3)$、
Hilbert級数$((1-t)^2(1-t^2)(1-t^3))^{-1}$を持つものを扱う。
全ての対象が量子平面またはJordan平面からの2段階の斉次Ore拡大であることを証明する。
最小解消の左右双対性により、重み3の関係の階数を重み4、5の関係へ伝えることが
網羅性の証明の核心となる。Oreデータの有限個の係数条件、
全ての重み保存同型を与える三角形変換、およびその必要十分条件を記述する。
さらに語長2次表示を持つ部分について、非対角自己同型を許す完全なgradedチャートを得る。
この結果の同値関係は次数付き代数の同型であり、一般のquadratic
$\ZZ$-代数の添字保持同型による全分類は別の問題として残る。
\end{abstract}
\vfill
\begin{center}
原稿作成日時：2026-11-18 14:22:06 JST（Asia/Tokyo）\\
組版ファイル：\texttt{\jobname.pdf}
\end{center}
\end{titlepage}
\hypersetup{pageanchor=true}
\pagenumbering{roman}
\tableofcontents
\clearpage
\pagenumbering{arabic}

\section{対象と分類定理}\label{sec:theorem}
基礎体を$k=\mathbb C$とする。代数と線型写像は全て$k$上で考える。
次数付き加群のshiftは$M(s)_n=M_{n+s}$とし、$A(-s)$の自由生成元の次数は$s$である。
代数の積にKoszul符号は入れない。
連結局所有限な$\mathbb N$-次数付き代数とは
$A=\bigoplus_{n\geq0}A_n$, $A_0=k$, $\dim_k A_n<\infty$を満たすものをいう。
$A_+=\bigoplus_{n>0}A_n$、$Q(A)=A_+/A_+^2$と置く。

\begin{definition}\label{def:graded-target}
$\Gclass$は次の条件を満たす代数の、次数を保つ同型類の集合とする。
\begin{enumerate}
\item $A$は連結局所有限で、左および右noetherianである。
\item 左右のglobal dimensionは4で、ある整数$\ell$について
\[
\Ext_A^i(k_A,A)=
\begin{cases}0&i\ne4,\\ {}_Ak(\ell)&i=4,\end{cases}
\qquad
\Ext_{A^{\mathrm{op}}}^i({}_Ak,A)=
\begin{cases}0&i\ne4,\\ k_A(\ell)&i=4\end{cases}
\]
というAS--Gorenstein条件を満たす。
\item 最小生成元空間は
$\dim Q(A)_1=2$, $\dim Q(A)_2=\dim Q(A)_3=1$、
他の次数では零である。
\item Hilbert級数は
\begin{equation}\label{eq:hilbert}
H_A(t)=\frac{1}{(1-t)^2(1-t^2)(1-t^3)}
\end{equation}
である。
\end{enumerate}
最後の条件は$\GKdim A=4$を含むので、有限成長を含む通常のAS正則性を満たす。
ここでは最小生成元次数だけから\eqref{eq:hilbert}を仮定なしに導くことは主張しない。
\end{definition}

$\sigma$を代数$B$の自己同型とする。
$\sigma$-導分とは
\[
\delta(ab)=\sigma(a)\delta(b)+\delta(a)b
\]
を満たす$k$-線型写像である。
Ore拡大$B[z;\sigma,\delta]$の規約は
$zb=\sigma(b)z+\delta(b)$とし、
左$B$-加群として$\bigoplus_{j\geq0}Bz^j$とする。
$|z|=d$なら$\sigma$の次数は0、$\delta$の次数は$d$である。

\begin{theorem}[Ore表示による網羅]\label{thm:ore-exhaustion}
$A$が$\Gclass$の対象であることは、次の表示を持つことと同値である。
\begin{equation}\label{eq:iterated-ore}
A=B[z;\sigma,\delta][w;\tau,D],\qquad |x|=|y|=1,\quad |z|=2,\quad |w|=3.
\end{equation}
ここで
\begin{equation}\label{eq:plane}
B=B_{q,\varepsilon}
=k\langle x,y\rangle/(r),\qquad r=yx-qxy-\varepsilon x^2,
\end{equation}
かつ
\[
(q,\varepsilon)\in
\{(q,0)\mid q\in k^\times\}\ \cup\ \{(1,1)\}.
\]
前者を量子平面（$q=1$では可換平面）、後者をJordan平面という。
$\sigma\in\Aut_{\mathrm{gr}}B$、
$\delta:B\to B$は次数2の$\sigma$-導分であり、
$C=B[z;\sigma,\delta]$と置いたとき
$\tau\in\Aut_{\mathrm{gr}}C$、
$D:C\to C$は次数3の$\tau$-導分である。

この表示では$x^iy^jz^hw^m$ ($i,j,h,m\geq0$) が$k$-基底となる。
$A$はdomainで、AS添字は7である。
中間部分代数は
\begin{equation}\label{eq:canonical-flag}
B=k\langle A_1\rangle,\qquad C=k\langle A_1,A_2\rangle
\end{equation}
として内在的に定まる。
\end{theorem}

定理の網羅性は第\ref{sec:exhaustion}節で証明する。
逆構成の正則性は第\ref{sec:converse}節で証明する。
第\ref{sec:data}節の有限係数条件と第\ref{sec:isomorphism}節の同型条件を合わせると、
明示したデータの同値類と$\Gclass$との全単射が得られる。
これは軌道による分類であり、全軌道の一意な標準代表やモジュライの成分分解を
追加して主張するものではない。

本稿の対象には語長3以上の項を含む関係も許す。
主原稿\path{nc-p1123.tex}は、各始点で最小生成元のjumpが$(1,1,2,3)$となる
語長2次の4次元AS正則$\ZZ$-代数を添字保持同型で分類する問題を扱い、
noetherian性、1周期性、所定のHilbert級数を全体には仮定していない。
両者の対応は第\ref{sec:quadratic}節で具体的に記述する。

\section{有限個の係数によるOreデータ}\label{sec:data}
\subsection{基礎平面の自己同型}
$V=kx\oplus ky$とする。次数付き自己同型群は次の通りである。
\begin{center}
\begin{tabular}{c|p{.65\linewidth}}
\toprule
$B$ & $\Aut_{\mathrm{gr}}B$の$V$への作用\\
\midrule
$B_{1,0}$ & 任意の$\GL(V)$の元\\
$B_{q,0}$, $q\ne\pm1$ & $x\mapsto ax,\ y\mapsto dy$, $a,d\ne0$\\
$B_{-1,0}$ & 可逆な対角行列または反対角行列\\
$B_{1,1}$ & $x\mapsto ax,\ y\mapsto bx+ay$, $a\ne0$\\
\bottomrule
\end{tabular}
\end{center}
実際、量子平面で$x'=ax+by$, $y'=cx+dy$と置くと
\[
y'x'-qx'y'=(1-q)ac\,x^2+(1-q^2)bc\,xy+(1-q)bd\,y^2.
\]
可逆性とこの3係数の消滅が表の量子平面の行を与える。
Jordan平面では$[y',x']-(x')^2$の$y^2$係数から$b=0$、
残る$x^2$係数と可逆性から$d=a$となる。
異なる基礎平面の間では、$B_{q,0}\simeq B_{q',0}$であることは
$q'=q$または$q'=q^{-1}$と同値であり、後者は$x,y$の交換による。
Jordan平面はこれらの量子平面とは同型でない。
この最後の区別は、関係テンソルの対称部分が定める二次形式の零点が、
非可換量子平面では2つの異なる直線、Jordan平面では重複直線となること、
および可換平面では対称部分が零であることから従う。
非可換量子平面の2直線を保つ、または交換する基底変更を代入すれば、
$q'=q^{\pm1}$という条件も得られる。

\subsection{導分と第2の自己同型}
平面$B=B_{q,\varepsilon}$と$\sigma\in\Aut_{\mathrm{gr}}B$を固定する。
次数2の$\sigma$-導分の候補を
\[
\delta(x)=f,\qquad \delta(y)=g,\qquad f,g\in B_3
\]
と書く。テンソル代数上では歪Leibniz則により一意に延長できる。
これが$B$に降りるための条件は、唯一の関係$r$の像が零であること、
すなわち
\begin{equation}\label{eq:delta-condition}
\sigma(y)f+gx-q\sigma(x)g-qfy
-\varepsilon\sigma(x)f-\varepsilon fx=0\quad\text{in }B_4
\end{equation}
である。$\dim B_3=4$, $\dim B_4=5$だから、
固定した$\sigma$の下では8係数に関する5本の線型方程式となる。
方程式の階数が常に5であるとは限らない。

$C=B[z;\sigma,\delta]$と置く。
任意の$\tau\in\Aut_{\mathrm{gr}}C$は、一意に
\begin{equation}\label{eq:tau-form}
\tau|_B=\mu,\qquad \tau(z)=cz+p,\qquad
\mu\in\Aut_{\mathrm{gr}}B,\quad c\in k^\times,\quad p\in B_2
\end{equation}
と書ける。この形の写像が自己同型となるための必要十分条件は
\begin{align}
\mu\sigma&=\sigma\mu,\label{eq:tau-commute}\\
c\,\delta(\mu(u))+p\mu(u)
-\mu(\sigma(u))p-\mu(\delta(u))&=0
\quad(u=x,y).\label{eq:tau-constant}
\end{align}
実際、$\tau$を$zu-\sigma(u)z-\delta(u)$に作用させ、
$z\mu(u)=\sigma(\mu(u))z+\delta(\mu(u))$を用いると、
$z$の係数が\eqref{eq:tau-commute}、定数項が\eqref{eq:tau-constant}になる。
これらが成立すれば準同型が定まり、その像は$B$を含み、
$c^{-1}(\tau(z)-p)=z$も含むので全射である。
各次数成分が有限次元であるから単射でもある。
逆に、任意の自己同型がこの形を持つことは
$B=k\langle C_1\rangle$、$C_2=B_2\oplus kz$から従う。

\subsection{最後の導分}
$C,\tau$を固定し、
\[
D(x)=F,\qquad D(y)=G,\qquad D(z)=H,\qquad
F,G\in C_4,\quad H\in C_5
\]
と置く。必要十分条件は
\begin{equation}\label{eq:D-plane}
\mu(y)F+Gx-q\mu(x)G-qFy
-\varepsilon\mu(x)F-\varepsilon Fx=0\quad\text{in }C_5
\end{equation}
と、$u=x,y$に対する
\begin{equation}\label{eq:D-mixed}
(cz+p)D(u)+Hu-\mu(\sigma(u))H
-D(\sigma(u))z-D(\delta(u))=0\quad\text{in }C_6
\end{equation}
である。右辺に現れる$D(\sigma(u))$、$D(\delta(u))$は、
指定した$F,G$と歪Leibniz則で計算する。
\eqref{eq:D-plane}は$D(r)=0$、
\eqref{eq:D-mixed}は$D(zu-\sigma(u)z-\delta(u))=0$そのものである。
関係イデアルはこれら3関係で生成されるため、この検査で降下の十分性も従う。

\begin{theorem}[有限係数条件]\label{thm:coefficient-data}
定理\ref{thm:ore-exhaustion}の全データは、
\eqref{eq:plane}の平面、上表の$\sigma,\mu$、$c\ne0$、
\[
f,g\in B_3,\quad p\in B_2,\quad F,G\in C_4,\quad H\in C_5
\]
を選び、\eqref{eq:delta-condition}、
\eqref{eq:tau-commute}--\eqref{eq:tau-constant}、
\eqref{eq:D-plane}--\eqref{eq:D-mixed}を課すことで尽くされる。
対応する代数の6関係は
\begin{equation}\label{eq:six-relations}
\begin{aligned}
yx&=qxy+\varepsilon x^2,&
zx&=\sigma(x)z+f,& zy&=\sigma(y)z+g,\\
wx&=\mu(x)w+F,&
wy&=\mu(y)w+G,& wz&=(cz+p)w+H.
\end{aligned}
\end{equation}
これらの条件を満たす全てのデータが$\Gclass$の対象を与える。
\end{theorem}
\begin{proof}
上の降下条件により$\delta,\tau,D$がそれぞれ定義される。
Ore拡大の普遍性から、その生成関係は\eqref{eq:six-relations}である。
逆にOreデータを生成元に制限すれば上の条件が全て必要となる。
正則性と他の対象条件は第\ref{sec:converse}節で確認する。
\end{proof}

順序基底$\{x^iy^jz^h\mid i+j+2h=n\}$で全式を展開できる。
\[
\dim C_4=9,\qquad \dim C_5=12,\qquad \dim C_6=16.
\]
従って最後の導分条件は、固定した$C,\tau$の下で
$9+9+12=30$係数を未知数とする
$(12+2\cdot16)\times30=44\times30$行列の核として記述できる。
この数は方程式の本数であり、独立条件の数や同型類空間の次元を表さない。

\section{全ての次数付き同型}\label{sec:isomorphism}
2つのデータから作る代数を
$A=B[z;\sigma,\delta][w;\tau,D]$、
$A'=B'[z';\sigma',\delta'][w';\tau',D']$とし、
$C=B[z;\sigma,\delta]$, $C'=B'[z';\sigma',\delta']$と置く。

\begin{theorem}[三角形変換による同型判定]\label{thm:graded-isomorphisms}
全ての次数付き同型$\Phi:A\to A'$は
\begin{equation}\label{eq:triangular}
\Phi|_B=\varphi,\qquad
\Phi(z)=\alpha z'+a_2,\qquad
\Phi(w)=\beta w'+b_3
\end{equation}
という形である。ただし
$\varphi:B\xrightarrow{\sim}B'$は次数付き同型、
$\alpha,\beta\in k^\times$、$a_2\in B'_2$、$b_3\in C'_3$である。
このデータが同型を与えるための必要十分条件は次の通りである。
まず
\begin{align}
\sigma'\varphi&=\varphi\sigma,\label{eq:iso-first-auto}\\
\varphi(\delta(u))
&=\alpha\delta'(\varphi(u))+a_2\varphi(u)
-\varphi(\sigma(u))a_2\quad(u=x,y).\label{eq:iso-first-der}
\end{align}
これらから誘導される同型を$\psi:C\to C'$とする。さらに
\begin{align}
\tau'\psi&=\psi\tau,\label{eq:iso-second-auto}\\
\psi(D(v))
&=\beta D'(\psi(v))+b_3\psi(v)-\psi(\tau(v))b_3
\quad(v=x,y,z).\label{eq:iso-second-der}
\end{align}
\eqref{eq:iso-second-auto}も$v=x,y,z$で検査すれば十分である。
\end{theorem}
\begin{proof}
\eqref{eq:canonical-flag}より任意の$\Phi$は$B,C$をそれぞれ$B',C'$へ移す。
$A_2=B_2\oplus kz$、$A_3=C_3\oplus kw$と最小生成元条件から
\eqref{eq:triangular}が得られ、$\alpha,\beta$は非零である。

$zu=\sigma(u)z+\delta(u)$を$\Phi$で移し、
$z'\varphi(u)=\sigma'(\varphi(u))z'+\delta'(\varphi(u))$と比較する。
$C'$の左$B'$-基底$1,z',\ldots$の係数を取ると
\eqref{eq:iso-first-auto}--\eqref{eq:iso-first-der}を得る。
逆にこの2式は$C\to C'$の準同型を定める。
その像が$B'$と$z'$を含むこと、各次数成分が有限次元で
$H_C=H_{C'}$であることから、この写像$\psi$は同型である。

続いて$wv=\tau(v)w+D(v)$を移すと、
左$C'$-基底$1,w',\ldots$の$w'$係数と定数項の比較により
\eqref{eq:iso-second-auto}--\eqref{eq:iso-second-der}を得る。
これらが生成元で成立すれば、自己同型と歪導分の積公式により
全ての$v\in C$で成立する。従ってOre拡大の普遍性から$\Phi$が定まり、
像が$C'$と$w'$を含むので全射である。
$H_A=H_{A'}$と次数成分の有限次元性から同型となる。
\end{proof}

以上の構成は、Oreデータから代数を作り、その内在的な列
$B\subset C\subset A$で生成元を選んでデータを取り出すという2操作を与える。
生成元の選択変更はちょうど定理\ref{thm:graded-isomorphisms}の変換である。
そのため、データをこの同型関係で割った集合から$\Gclass$への写像は、
網羅定理により全射、同型判定により単射である。

\section{網羅性の証明}\label{sec:exhaustion}
最小解消の双対性から関係の個数と次数を定め、
最小重み4のsyzygyの係数を高い重みの関係へ移す。
まず$x,y,z$の関係をOre型に強制し、その階数情報から$w$もOre型にする。
最後に所定のHilbert級数を使って、得られた順序単項式の独立性を示す。

\subsection{domain性と最小解消}
$A\in\Gclass$を固定する。
\eqref{eq:hilbert}から$\GKdim A=4$である。
noetherian連結次数付きAS正則代数について、global dimensionとGK次元が
ともに4ならdomainであるという
Lu--Palmieri--Wu--Zhang\cite[Lemma~1.3]{LPWZ}を適用する。
同補題は次数1生成を仮定しないので、今回の最小生成元次数にも適用できる。

最小右自由解消を$P_\bullet\to k_A$とする。
AS--Gorenstein双対性により
$\Hom_A(P_\bullet,A)(-\ell)$は、順序を逆にして左$A$-加群$k$の最小自由解消となる。
左右の最小解消の自由生成元次数は同じ$\Tor^A_i(k,k)$で与えられるので、
次数の多重集合はホモロジー次数$i$と$4-i$で$\ell$に関して対称である
\cite[Lemma~1.1]{LPWZ}。
Euler多項式$p(t)=H_A(t)^{-1}$は従って
$p(t)=t^\ell p(t^{-1})$を満たす。
$p(t)=(1-t)^2(1-t^2)(1-t^3)$の定数項は1、次数は7なので$\ell=7$である。
$P_1$は最小生成元条件から既知であり、双対性で$P_3,P_4$を定めた後に
Euler等式で$P_2$を定めると、次を得る。
\begin{equation}\label{eq:resolution}
\begin{aligned}
P_0&=A,& P_1&=A(-1)^2\oplus A(-2)\oplus A(-3),\\
P_2&=A(-2)\oplus A(-3)^2\oplus A(-4)^2\oplus A(-5),\\
P_3&=A(-4)\oplus A(-5)\oplus A(-6)^2,& P_4&=A(-7).
\end{aligned}
\end{equation}
従って最小関係は重み$2,3,3,4,4,5$の6本である。
以下で「二次・三次」等と呼ぶ次数は、この節では重み付き次数を指す。
最小生成元を$x,y,z,w$と選ぶと、二次関係に$z$の線型項、
三次関係に$w$の線型項は現れない。

\subsection{次数1部分が生成する平面}
唯一の二次関係$r\in V\otimes V$、$V=kx\oplus ky$は行列階数2である。
階数1なら2つの非零な次数1の元の積が$A$で零となり、domain性に反する。
この非退化な二変数の関係は、基底変更と関係の非零定数倍により
\eqref{eq:plane}の形になる。具体的には、関係行列の対称部分の二次形式が
零なら交代形式、2つの異なる零方向を持てば両対角成分を零にした行列、
重複した零方向だけを持てば
$\left(\begin{smallmatrix}a&b\\-b&0\end{smallmatrix}\right)$
($ab\ne0$)にできる。
それぞれ可換平面、量子平面、Jordan平面の関係となる。
行列の非退化性により量子パラメータ$q$は非零である。

\begin{lemma}\label{lem:plane-embedding}
$B=B_{q,\varepsilon}\to A$という自然な準同型は単射である。
\end{lemma}
\begin{proof}
$B$では$yx=qxy+\varepsilon x^2$という規則に自己重なりはなく、
この規則はdegree-lex順$y>x$を下げる。
従って$x^iy^j$が基底となる。
$A$はnoetherian domainなので、Goldieの定理により分数斜体$K$に埋め込める。
$t=yx^{-1}\in K$と置く。
量子平面の場合は$tx=qxt$であり、
\[
x^{n-j}y^j=q^{j(j+1)/2}x^n t^j.
\]
Jordan平面の場合は$tx=x(t+1)$であり、
\[
x^{n-j}y^j=x^n\prod_{h=1}^{j}(t+h).
\]
いずれの場合も、$B_n$の非零元の像は$x^nP(t)$と書け、
$P\in k[T]$は非零である。写像の核は斉次だから、
単射でなければある非零の$P$について$P(t)=0$となる。
$k=\mathbb C$上で$P$を一次式に分解し、$K$が斜体であることを使うと、
ある$\lambda\in k$について$t=\lambda$となる。
これは$y=\lambda x$を意味し、$\dim A_1=2$に反する。
\end{proof}

\subsection{双対解消が伝える係数}\label{subsec:dual-coefficient}
自由代数を$T=k\langle x,y,z,w\rangle$、関係イデアルを$I$とする。
最小関係を$\rho$で添字付け、$b_\rho=\deg r_\rho$と置く。
右解消の$P_2$の基底を$e_\rho$とし、
$r_\rho=\sum_{v=x,y,z,w}v f_{v\rho}$という左端の文字による分解から
$d_2(e_\rho)=\sum_v e_v f_{v\rho}$を取る。
$P_3$の基底の次数は$s=4,5,6,6$であり、次数6の2本は別々に添字付ける。
\[
d_3(e_s)=\sum_\rho e_\rho N_{\rho s},\qquad
d_4(e)=\sum_s e_s q_s.
\]
ここで$\deg N_{\rho s}=s-b_\rho$、$\deg q_s=7-s$である。
双対左解消の最初の写像は$q_s$を生成元とし、その最小関係は
\begin{equation}\label{eq:dual-relations}
r^\vee_\rho=\sum_s \widetilde N_{\rho s}\,\widetilde q_s
\end{equation}
で与えられる。チルダは$T$への斉次持上げを表す。
実際、$d_3d_4=0$がこの式を$A$内で零にし、
双対左解消の正確性と最小性が、これらを最小関係の基底とする。
左端の写像が$A_+$を最小生成するので、
\begin{equation}\label{eq:dual-generators}
\begin{gathered}
(q_{6,1},q_{6,2})\text{ は }V\text{ の基底},\\
q_5=\alpha z+p_2,\qquad q_4=\beta w+p_3,\qquad
\alpha\beta\ne0,\quad
p_2\in B_2,\quad p_3\in B_3+Vz+zV.
\end{gathered}
\end{equation}
この生成元変更は自由代数でも三角形の可逆変換である。

次数4の$P_3$生成元を$e_4$とし、その$r_2$行と2本の$r_3$行を
\begin{equation}\label{eq:lowest-syzygy}
N_{2,4}=\gamma z+p_2',\qquad
N_{3,1;4}=b_1,\qquad N_{3,2;4}=b_2,\qquad
p_2'\in B_2,\quad b_1,b_2\in V
\end{equation}
と書く。最小性により、同次数4の関係の行には定数係数はなく、他の行も零である。

\begin{lemma}[低次syzygyと高次関係の係数]\label{lem:dual-coefficients}
\eqref{eq:lowest-syzygy}において、
\begin{enumerate}
\item $\gamma=0$であることと、唯一の最小五次関係の$zw$係数が零であることは同値である。
\item $b_1,b_2$が$V$の基底であることと、2本の最小四次関係から$Vw$への
係数射影が階数2であることは同値である。
\end{enumerate}
反対代数に対して同じ記述を適用すると、$wz$係数と$wV$への射影を得る。
\end{lemma}
\begin{proof}
\eqref{eq:dual-relations}の$r_2$に対応する行は次数5である。
その$zw$項は$s=4$の
$(\gamma z+p_2')(\beta w+p_3)$からだけ生じ、係数は$\beta\gamma$である。
$s=5$では右因子が次数2で、$s=6$では右因子が$V$に入るため、
他の項から$zw$は生じない。
また、2本の$r_3$行は次数4の双対関係であり、
その$Vw$部分は$s=4$の$\beta b_1w,\beta b_2w$だけである。

これらの双対関係を元の最小関係と比較する。
四次関係は同次数の2本の可逆な線型結合と低次関係の倍だけ異なる。
低次関係は$w$を含まず、それらに掛ける係数の次数は高々2なので、
$Vw,wV$係数はその可逆な線型結合だけを受ける。
五次関係は非零定数倍と低次関係の倍だけ異なる。
$(V)$を法とすると二次・三次関係は零であり、
四次関係の倍にも必ず次数1の因子が付くので零になる。
従って$zw,wz$係数は非零定数倍だけを受ける。
持上げの変更も同じ低次関係の倍に含まれる。
\eqref{eq:dual-generators}で$\beta\ne0$だから、2つの主張が従う。
\end{proof}

\subsection{三次関係の片側射影}
二次関係を先に割ると、自由代数の次数3部分は
\[
B_3\oplus Vz\oplus zV\oplus kw
\]
と分解される。2本の最小三次関係には$w$項がない。
それらの$Vz\oplus zV$への射影の像を$R$とする。
補題\ref{lem:plane-embedding}により、三次関係の非零な線型結合が
$B_3$だけに入ることはないので$\dim R=2$である。

\begin{lemma}\label{lem:cubic-projection}
$R$から$Vz$、$zV$への2つの射影の少なくとも一方は階数2である。
\end{lemma}
\begin{proof}
どちらも階数2でないとする。片方の階数が零なら他方は2となるので、
両方の階数は1である。両射影の核は異なる直線であり、
それぞれの核の基底を取ると三次関係を
\[
r_{3,1}=az-f,\qquad r_{3,2}=zb-h,\qquad
0\ne a,b\in V,\quad f,h\in B_3
\]
と書ける。
$r_2=xu_x+yu_y$とすると、関係行列の非退化性により$u_x,u_y$は$V$の基底である。
次数4の右syzygyの係数を
$(\gamma z+p,\ell_1,\ell_2)$と書く。
ここで$p\in B_2$, $\ell_1,\ell_2\in V$であり、残る関係の係数は最小性により零である。
$d_2d_3=0$の$z$行は$b\ell_2=0$であるから、domain性により$\ell_2=0$となる。
$a=a_xx+a_yy$と書き、残る$x,y$行を$B_3$で割ると
\[
\gamma u_jz+a_jz\ell_1=0
\quad\text{in }(Vz\oplus zV)/R,\qquad j=x,y.
\]
もし$\gamma\ne0$なら、$Vz$への射影を取ると
$u_x,u_y\in ka$となり、独立性に反する。従って$\gamma=0$である。

反対代数でも両射影の階数は1であり、同じ計算で対応する係数が零となる。
補題\ref{lem:dual-coefficients}から、最小五次関係の$zw,wz$係数はいずれも零である。
よって$(x,y)$を法とした表示では、二次・三次関係は零、
四次関係は$z^2$の定数倍、五次関係は零となる。
従って$A/(x,y)$は
$k\langle z,w\rangle$または$k\langle z,w\rangle/(z^2)$である。
後者でも$w,zw$の任意の語は$z^2$を含まず、相異なる語は相異なる単項式となる。
従って両者は指数増大する。これは$A$の有限GK次元に矛盾する。
\end{proof}

\subsection{最初のOre自己同型の可逆性}
まず$R\to zV$の射影が可逆である場合を考える。
$x_1=x$, $x_2=y$とすると、三次関係は
\begin{equation}\label{eq:cubic-graph}
r_{3,i}=zx_i-\sigma(x_i)z-\delta_i,\qquad \delta_i\in B_3
\end{equation}
と書ける。この段階では$\sigma:V\to V$は線型写像である。
$X=(x_1,x_2)^T$、
$r_2=X^TMX$、$\sigma(X)=SX$と置く。
$M$は可逆である。
最低次数の右syzygyを\eqref{eq:lowest-syzygy}のように書くと、
$d_2d_3=0$の$z$行は
\[
x_1b_1+x_2b_2=0.
\]
二次関係が$r_2$だけであることから、
\begin{equation}\label{eq:b-vector}
(b_1,b_2)^T=\lambda MX
\end{equation}
となる。
ここで$\lambda\ne0$である。
実際、$\lambda=0$ならsyzygyは$r_2$行だけとなり、
$(MX)_j(\gamma z+p_2')=0$とdomain性からこの行も零になる。
しかし$d_3(e_4)=0$なら$e_4\in\operatorname{im}d_4$となり、
$P_4$の最小次数が7であることに反する。

$x_j$行の式は
\[
(MX)_j(\gamma z+p_2')
-\sum_{i=1}^2 S_{ij}zb_i
-\sum_{i=1}^2\delta_{ij}b_i=0,
\qquad
\delta_i=\sum_{j=1}^2x_j\delta_{ij}.
\]
$\delta_{ij}\in B_2$であり、
\eqref{eq:cubic-graph}を使って$zb_i$を置き換える。
$Vz$は$B_3$を法として独立であるから、
$Vz$係数を比較し\eqref{eq:b-vector}を用いると
\begin{equation}\label{eq:matrix-identity}
\gamma M=\lambda S^TMS
\end{equation}
を得る。従って
$\sigma(r_2)=(\gamma/\lambda)r_2$というテンソル代数内の等式が成立し、
$\sigma$は$B$の自己準同型に降りる。
$\delta(x_i)=\delta_i$を歪Leibniz則で延長すると、
$zr_2=0$と\eqref{eq:cubic-graph}から$\delta(r_2)=0$が$A$内で成立する。
$\delta(r_2)\in B_4$と補題\ref{lem:plane-embedding}により、
これは$B$内でも零であり、$\delta$は$B$の$\sigma$-導分となる。

\begin{lemma}\label{lem:first-automorphism}
$S$は可逆である。従って$\sigma\in\Aut_{\mathrm{gr}}B$、
$\gamma\ne0$であり、$b_1,b_2$は$V$の基底である。
\end{lemma}
\begin{proof}
$S$が特異なら、$\ker\sigma$は$B$の非零な次数1のnormal元$v$を含む。
これを各平面で確認する。
可換平面では任意の次数1の元がnormalである。
非可換量子平面で$\rk S=1$なら
$\sigma(x)=a u$, $\sigma(y)=b u$と書け、
関係の像は$(1-q)ab\,u^2=0$であるから$ab=0$となる。
従って$x$または$y$が核に入り、いずれもnormalである。
Jordan平面で$\rk S=1$なら同じ表示から
$-a^2u^2=0$、従って$\sigma(x)=0$となり、normal元$x$が核に入る。
$S=0$の場合も各平面で$x$を選べばよい。

$zv=\delta(v)=f\in B_3$とする。
$v$のnormal性から$vf=hv$となる$h\in B_3$が存在する。
従って$(vz-h)v=0$となり、$A$のdomain性から$vz=h$を得る。
これは$zV$成分が零で$Vz$成分が非零の三次関係であり、
\eqref{eq:cubic-graph}が$zV$上のグラフであることに反する。
よって$S$は可逆である。
\eqref{eq:matrix-identity}から$\gamma\ne0$、
\eqref{eq:b-vector}から$b_1,b_2$の独立性が従う。
また$\sigma(r_2)$が$r_2$の非零定数倍なので、
$S^{-1}$も関係を保ち、$\sigma$は自己同型である。
\end{proof}

補題\ref{lem:cubic-projection}で$Vz$側しか可逆と分からなかった場合は、
直前の議論を$A^{\mathrm{op}}$に適用する。
そこで得る線型写像も可逆であるため、元の$A$に戻って
$zV$についても関係を解ける。このように、いずれの場合も$A$自身について
\eqref{eq:cubic-graph}を可逆な$\sigma$で取れる。
さらに逆向きの関係
\[
uz=z\sigma^{-1}(u)-\delta(\sigma^{-1}(u))
\]
があるので、反対代数にも補題\ref{lem:first-automorphism}を適用できる。
補題\ref{lem:dual-coefficients}により、
四次関係の$wV,Vw$への射影はいずれも階数2であり、
五次関係の$wz,zw$係数はいずれも非零である。

\subsection{最後のOre付加とHilbert級数}
四次関係の$wV$への射影が可逆であることから、
2本の四次関係を
\[
wx=\mu(x)w+F,\qquad wy=\mu(y)w+G
\]
と書ける。ここで$\mu:V\to V$は線型写像、
$F,G$は$x,y,z$だけを含む次数4の式である。
$Vw$への射影も可逆なので、$\mu|_V$は可逆である。
五次関係では、これらの四次関係を用いて
$wV^2,VwV$を右端に$w$を持つ項に置き換える。
$wz,zw$係数はいずれも非零なので
\[
wz=(cz+p)w+H,\qquad c\ne0,\quad p\in B_2
\]
とできる。$H$は$x,y,z$だけを含む次数5の式である。
この置換では$zw,wz$係数は変わらない。

ここで形式的なOre拡大$C_0=B[z;\sigma,\delta]$を作る。
まだ$C_0\to A$の単射性は使わない。
得られた6関係を$yx,zx,zy,wx,wy,wz$の書換え規則とする。
重み付き次数を先に比較し、同次数では$w>z>y>x$による辞書式順序を使うと、
全ての右辺の語は左辺より小さい。正の重みなので各次数の語は有限個であり、
同重みの異なる語が互いの真の接頭語となることはないため、
この辞書式順序は左右の乗法と両立し、書換えが停止する。
従って$x^iy^jz^hw^m$が$A$を張る。
これらの重み$n$の個数は\eqref{eq:hilbert}の$t^n$係数に等しいため、
各次数で線型独立である。これにより
\begin{equation}\label{eq:PBW}
C_0\hookrightarrow A,\qquad A=\bigoplus_{m\geq0}C_0w^m
\end{equation}
を得る。以後$C=C_0$と同定する。

$C$の生成元との積を上の規則で書き換えると、
任意の$a\in C$について一意に
$wa=\tau(a)w+D(a)$と書ける。
$\tau(1)=1$, $D(1)=0$であり、$w(ab)=(wa)b$と
\eqref{eq:PBW}の係数比較により
\[
\tau(ab)=\tau(a)\tau(b),\qquad
D(ab)=\tau(a)D(b)+D(a)b
\]
となる。従って$\tau$は次数0の自己準同型、$D$は次数3の$\tau$-導分である。
$\tau|_V=\mu$は可逆で、$\tau(z)=cz+p$, $c\ne0$だから、
$\tau(C)$は$B$と$z$を含む。よって$\tau$は全射であり、
局所有限性と次数保存から自己同型である。
Ore拡大$C[w;\tau,D]\to A$は\eqref{eq:PBW}の基底を保つので同型である。
$B=k\langle A_1\rangle$、$C=k\langle A_1,A_2\rangle$も基底から従い、
定理\ref{thm:ore-exhaustion}の網羅性が証明された。

\section{逆構成の正則性}\label{sec:converse}
平面\eqref{eq:plane}は
$k[x][y;\nu,\eta]$と書ける。
ここで$\nu(x)=qx$、$\eta(x)=\varepsilon x^2$であり、
$\nu$は自己同型、$\eta$は次数1の$\nu$-導分である。
従って$B,C,A$は、次数1の多項式環$k[x]$から
自己同型を使って順次Ore付加した代数である。

自己同型によるOre拡大は左右noetherian性とdomain性を保つ
\cite[\S1.3]{LWW}。
またLiu--Wang--Wu\cite[Theorem~3.3]{LWW}によれば、
次元$d$のtwisted Calabi--Yau代数へのこのOre付加は
次元$d+1$のtwisted Calabi--Yau代数となる。
$k[x]$は次元1のtwisted Calabi--Yau代数であるから、
順に$B,C,A$の次元は2、3、4となる。
同論文Lemma~1.3により、次数を忘れた自明加群のExtは最上次数にだけ存在し、
その次元は1である。これを有限生成斉次自由加群による解消で計算すると、
次数を忘れたHom複体は内部次数ごとのHom複体の直和となる。
従ってその1次元Extはある内部次数に集中し、次数付き加群として$k(\ell)$である。
左右の最上Torの次数は一致するため、左右のAS添字も一致する。
これで定義\ref{def:graded-target}の次数付きAS--Gorenstein条件を得る。
正次数のOre変数を持つ場合の基本的な双加群短完全列のshiftは、
同論文Remark~2.2にも明記されている。
同論文の定義には有限GK成長を含めないので、
ここでは次の基底計算でその条件を別に確認する。

各Ore付加は左係数環上の自由加群であるから
$x^iy^jz^hw^m$が基底であり、Hilbert級数は\eqref{eq:hilbert}となる。
これにより$\GKdim A=4$である。
全ての生成関係は$x,y,z,w$の2文字以上の語からなるので、
関係イデアルは自由代数のaugmentation idealの2乗に含まれ、
$Q(A)$の基底は$x,y,z,w$の像である。
したがって最小生成元次数も指定通りである。
最後にAS双対性とこのHilbert級数に対する
第\ref{sec:exhaustion}節冒頭のEuler計算でAS添字7を得る。
これで定理\ref{thm:ore-exhaustion}の逆向きが証明された。

\section{語長2次部分と既存のOreチャート}\label{sec:quadratic}
語長とは、選んだ最小生成元$x,y,z,w$の各文字を1と数える次数をいう。
語長2次表示を持つことは、ある斉次持上げの選択で
関係イデアルを語長2の関係で生成できることを意味する。

\begin{corollary}[非対角自己同型を許す二次チャート]\label{cor:quadratic-chart}
$\Gclass$のうち語長2次表示を持つ代数は、
その表示に適合する生成元で次の形に書ける。
\begin{equation}\label{eq:quadratic-relations}
\begin{gathered}
B=B_{q,\varepsilon},\qquad
A=B[z;\sigma][w;\tau,D],\\
\sigma,\mu\in\Aut_{\mathrm{gr}}B,\qquad
\tau|_B=\mu,\quad \tau(z)=cz,\quad c\in k^\times,\\
D(u)=L(u)z^2\quad(u\in V),\qquad D(z)=0,\quad L\in V^*.
\end{gathered}
\end{equation}
必要十分条件は
\begin{equation}\label{eq:quadratic-conditions}
\begin{gathered}
\mu\sigma=\sigma\mu,\qquad L\circ\sigma=cL,\\
L(x)\bigl(\mu(y)-q\sigma^2(y)
-\varepsilon\mu(x)-\varepsilon\sigma^2(x)\bigr)
+L(y)\bigl(\sigma^2(x)-q\mu(x)\bigr)=0
\quad\text{in }V.
\end{gathered}
\end{equation}
逆にこの条件を満たす全データが、語長2次表示を持つ$\Gclass$の対象を与える。
\end{corollary}
\begin{proof}
語長2次表示を持つ最小生成元を初めに選ぶ。
網羅性の証明では$x,y$の線型変更と関係の基底変更だけを使い、
$z,w$を非線型に変更していない。
重み3、語長2の語は$Vz,zV$だけなので、
最初のOre表示では$\delta=0$である。
重み4、語長2の語は$Vw,wV,z^2$だけなので
$D(u)=L(u)z^2$となる。
重み5、語長2の語は$zw,wz$だけなので
$\tau(z)=cz$、$D(z)=0$となる。
従って\eqref{eq:quadratic-relations}が必要である。

$zu=\sigma(u)z$を$\tau$で保つ条件は$\mu\sigma=\sigma\mu$である。
歪導分の降下条件は
\[
D(zu-\sigma(u)z)=
\bigl(cL(u)-L(\sigma(u))\bigr)z^3
\]
および
\[
\begin{split}
D(r)=\bigl[
&L(x)\bigl(\mu(y)-q\sigma^2(y)
-\varepsilon\mu(x)-\varepsilon\sigma^2(x)\bigr)\\
&+L(y)\bigl(\sigma^2(x)-q\mu(x)\bigr)\bigr]z^2
\end{split}
\]
が零となることである。$C=B[z;\sigma]$の基底により、
これは\eqref{eq:quadratic-conditions}と同値である。
逆向きは定理\ref{thm:coefficient-data}と第\ref{sec:converse}節による。
\end{proof}

この系の$\delta=0$等は、語長2次表示に適合した座標に関する条件である。
一般の三角形変換後も同じ係数が零とは限らない。
第\ref{sec:data}節の全データ空間の中で、
語長2次表示を持つ代数の部分は、上のチャートの
定理\ref{thm:graded-isomorphisms}の同型関係による飽和として得られる。

\begin{proposition}[二次部分での同型の存在]\label{prop:quadratic-isomorphisms}
2つの二次データ$(B,\sigma,\mu,c,L)$、
$(B',\sigma',\mu',c',L')$が定める代数が次数付き同型となることは、
次数付き同型$\varphi:B\to B'$と$\alpha,\beta\in k^\times$が存在して
\begin{equation}\label{eq:quadratic-isomorphism}
\sigma'\varphi=\varphi\sigma,\qquad
\mu'\varphi=\varphi\mu,\qquad
c'=c,\qquad
\beta L'\circ\varphi=\alpha^2L
\end{equation}
となることと同値である。
最後の式は$V$上の線型形式の等式である。
\end{proposition}
\begin{proof}
条件が成立すれば
$u\mapsto\varphi(u)$ ($u=x,y$)、$z\mapsto\alpha z'$、
$w\mapsto\beta w'$が6関係を保ち、同型を与える。
例えば$wu=\mu(u)w+L(u)z^2$を移すと
$z'^2$係数の比較が最後の条件を与え、
$wz=czw$を移すと$c'=c$となる。

逆に任意の次数付き同型はaugmentation ideal $I_A=A_+$の冪を保つ。
語長$n$の成分を$A^{(n)}$と書くと、語長斉次表示では
$I_A^m=\bigoplus_{n\geq m}A^{(n)}$である。
各重みでは語長が有限範囲にあるため、
選んだ語長次数により$\gr_{I_A}A\simeq A$となる。
同型に付随次数を取れば、重みと語長をともに保つ同型を得る。
その語長1の作用は必ず
$u\mapsto\varphi(u)$、$z\mapsto\alpha z'$、
$w\mapsto\beta w'$であり、関係を比較すると
\eqref{eq:quadratic-isomorphism}が必要となる。
\end{proof}

この命題は同型の存在を線型データで判定する。
非線型項を含む全ての同型そのものは、引き続き
定理\ref{thm:graded-isomorphisms}で与えられる。

\subsection{対角チャートへの特殊化}
\path{papers/weighted_ore_strata.tex}の記号との対応は
\[
\begin{gathered}
\varepsilon=0,\quad \sigma(x)=px,\quad \sigma(y)=ry,\quad
\mu(x)=sx,\quad \mu(y)=uy,\\
c=v,\qquad L(x)=a,\qquad L(y)=b,\qquad w=t
\end{gathered}
\]
である。このとき\eqref{eq:quadratic-conditions}は
\[
a(v-p)=a(u-qr^2)=b(v-r)=b(p^2-qs)=0
\]
となり、同原稿の4係数条件を回収する。
この対角表示空間の成分分解は、同原稿の\texttt{thm:components}に記述されている。
$q\ne\pm1$の量子平面では$\Aut_{\mathrm{gr}}B$が対角なので、
対角チャートがこの基礎平面上の二次graded部分を尽くす。
可換平面、$q=-1$の量子平面、Jordan平面では上記の非対角自己同型も必要になる。

\subsection{添字保持同型との対応}
次数付き代数$A$に付随する$\ZZ$-代数は
$\check A_{ij}=A_{j-i}$ ($j\geq i$)、$i>j$では零とし、
積を$A$の積で定める。
語長2次部分の$\check A$は、各始点で最小生成元のjumpが$(1,1,2,3)$である
語長2次AS正則$\ZZ$-代数で、標準1周期構造を持つ。
実際、$A$の最小自由解消の各$A(-d)$を始点$i$で$e_{i+d}\check A$に置き換えると、
各右単純加群の解消を得る。左解消とAS双対性も同じ添字変更で得られ、
各行のHilbert級数は\eqref{eq:hilbert}である。
次数付き$A$-加群と右$\check A$-加群は各次数成分を各頂点成分とする対応で同一視され、
この対応は射影加群と解消を保つ。斉次全射の次数を忘れた分裂から次数0成分を取れるので、
$A$のglobal dimension 4は次数付き解消の長さの上界でもある。
上記の単純加群の最小解消の長さが4なので、$\check A$のglobal dimensionも4である。
次数付き同型はこの標準1周期構造を保つ添字保持同型に対応する。
標準周期構造を忘れた添字保持同型ではさらにZhang twistが同一視される
\cite[Theorem~1.2]{Sierra}。
従って\eqref{eq:quadratic-isomorphism}を、そのまま裸の$\check A$の同型判定には使わない。
本稿は一般の7周期的対象やnoetherian性を仮定しない対象を分類していない。
また、一般の語長2次AS正則$\ZZ$-代数において所定のHilbert級数が
強制されるかという全体被覆の問題も残る。

\section{具体例と中心生成元}
\subsection{対角量子多項式環}
$q,a,b,c,d,e\in k^\times$を任意に取り、
\[
yx=qxy,\quad zx=axz,\quad zy=byz,\quad
wx=cxw,\quad wy=dyw,\quad wz=ezw
\]
と置けば$\Gclass$の対象を得る。
データでは$\delta=D=0$、
$\sigma=\operatorname{diag}(a,b)$、
$\mu=\operatorname{diag}(c,d)$である。
2つの対角自己同型は可換であるため、係数条件は全て成立する。

\subsection{可換三生成元環への導分付加}
$C=k[x,y,z]$、$\deg(x,y,z)=(1,1,2)$、$\tau=\mathrm{id}$とする。
任意の$F,G\in C_4$, $H\in C_5$に対し
\[
D=F\frac{\partial}{\partial x}
 +G\frac{\partial}{\partial y}
 +H\frac{\partial}{\partial z}
\]
は次数3の導分であり、$C[w;D]\in\Gclass$である。
この表示の係数空間は$9+9+12=30$次元である。
同型類は定理\ref{thm:graded-isomorphisms}でさらに同一視されるので、
30を同型類空間の次元としない。

\begin{corollary}[最高次数の最小生成元が中心の場合]\label{cor:central-generator}
$A\in\Gclass$が、$Q(A)_3$に非零の像を持つ中心斉次元$w_0\in A_3$を持つなら、
\[
A=C[w_0],\qquad C=B[z;\sigma,\delta]
\]
という通常の多項式拡大である。
逆に任意のこのような$C$について$C[w_0]\in\Gclass$であり、$w_0$は中心である。
\end{corollary}
\begin{proof}
網羅定理の$w$に対し$w_0=\beta w+b_3$、
$\beta\ne0$, $b_3\in C_3$と書ける。
各$w_0^m$の$w$次数の最高項は$\beta^m w^m$であるから、
$A=\bigoplus_{m\geq0}Cw_0^m$でもある。
$w_0a=aw_0$ ($a\in C$)のこの基底での係数比較により、
$w_0$によるOre自己同型は恒等写像、導分は零となる。
従って$C[w_0]\to A$は基底を保つ同型である。逆は逆構成から従う。
\end{proof}
\begingroup
\small
\begin{thebibliography}{9}
\bibitem{LPWZ}
D.-M.~Lu, J.~H.~Palmieri, Q.-S.~Wu and J.~J.~Zhang,
\emph{Regular algebras of dimension 4 and their $A_\infty$-Ext-algebras},
Duke Mathematical Journal \textbf{137} (2007), 537--584.
\href{https://doi.org/10.1215/S0012-7094-07-13734-7}{doi:10.1215/S0012-7094-07-13734-7}.
著者版：\url{https://sites.math.washington.edu/~palmieri/Papers/hiext.pdf}.
使用箇所：Lemma~1.1（左右最小解消の双対性）、
Lemma~1.3（次数1生成を要求しないdomain定理）。
\bibitem{LWW}
L.~Liu, S.~Wang and Q.~Wu,
\emph{Twisted Calabi--Yau property of Ore extensions},
Journal of Noncommutative Geometry \textbf{8} (2014), 587--609.
\href{https://doi.org/10.4171/JNCG/165}{doi:10.4171/JNCG/165}.
出版版：\url{https://ems.press/content/serial-article-files/30578}.
使用箇所：\S1.2、Lemma~1.3、\S1.3、Remark~2.2、Theorem~3.3。
\bibitem{Sierra}
S.~J.~Sierra,
\emph{$G$-algebras, twistings, and equivalences of graded categories},
Algebras and Representation Theory \textbf{14} (2011), 377--390.
\href{https://doi.org/10.1007/s10468-009-9193-y}{doi:10.1007/s10468-009-9193-y};
\href{https://arxiv.org/abs/math/0608791}{arXiv:math/0608791}.
使用箇所：Theorem~1.2（次数1生成を仮定しないtwistと付随$\ZZ$-代数の対応）。
\end{thebibliography}
\endgroup
\end{document}
